\subsection{同态的伴随。}

本节中，我们用 A 表示一个环（相应地一个配备了反自同构 J 的环），用 E 与 E′ 表示两个左 A-模，用 F 与 F′ 表示两个右 A-模（相应地左 A-模），并用 Φ 与 Φ′ 分别表示 E × F 与 E′ × F′ 上的两个双线性型（相应地关于 J 的半双线性型）。我们假设 Φ 非退化，换言之（第 1 小节），与 Φ 相伴的线性映射 \(d_{\Phi}\) 与 \(s_{\Phi}\) 都是单射。

给定一个同态 \(u\) （从 E 到 \(\mathbf{E}'\) 中），考虑这样的集合 \(\mathbf{F}_{1}'\) ：它的元素 \(y'\) 属于 \(\mathbf{F}'\) ，且存在 \(y \in \mathbf{F}\) 使得有 \(d_{\Phi'}(y') \circ u = d_{\Phi}(y)\) ，也就是说 \(\Phi'(u(x), y') = \Phi(x, y)\) （对一切 \(x \in \mathbf{E}\) ）。显然 \(\mathbf{F}_{1}'\) 是 \(\mathbf{F}'\) 的一个子模。由于 \(d_{\Phi}\) 是单射，对每个 \(y' \in \mathbf{F}_{1}'\) ，存在且仅存在 F 的一个元素 \(y\) ，使得 \(\Phi'(u(x), y') = \Phi(x, y)\) 。如此定义的映射 \(y' \to y\) （从 \(\mathbf{F}_{1}'\) 到 F 中）是 A-线性的；把它记作 \(u^{*}\) ，则对每个 \(x \in \mathbf{E}\) 与每个 \(y' \in \mathbf{F}_{1}'\) ，有

\[
\Phi^ {\prime} (u (x), y ^ {\prime}) = \Phi (x, u ^ {*} (y ^ {\prime})).\tag{30}
\]

定义 10. --- 在上述假设与记号下，称同态 \(u^*\) （从 \(F_1'\) 到 F 中，满足 (30)）为 \(u\) 的左伴随，而称 \(F_1'\) 为 \(u^*\) 的定义子模。

类似地，用公式定义 F 到 F′ 中的同态 \(\nu\) 的右伴随

\[
\Phi^ {\prime} (x ^ {\prime}, \nu (y)) = \Phi (\nu^ {*} (x ^ {\prime}), y) \quad (x ^ {\prime} \in \mathrm{E} _ {1} ^ {\prime}, y \in \mathrm{F}),\tag{31}
\]

其中 \(E_{1}^{\prime}\) 表示 \(E^{\prime}\) 的、与 \(F_{1}^{\prime}\) 类似地定义的子模。

注记. --- 如果 \(u^{*}\) （\(u: E \to E'\) 的左伴随）处处有定义，且 \(s_{\Phi'}\) 与 \(d_{\Phi'}\) 都是单射，则公式 (30) 表明 \(u\) 是 \(u^{*}\) 的右伴随。

由 (30) 我们导出：若 \(u_1\) 与 \(u_2\) 是 E 到 E′ 中的两个同态，其伴随都处处有定义，且 c 是 A 的中心的一个元素，则有

\[
\left\{ \begin{array}{l} (u _ {1} + u _ {2}) ^ {*} = u _ {1} ^ {*} + u _ {2} ^ {*}; 1 ^ {*} = 1; \\ (c u _ {1}) ^ {*} = c. u _ {1} ^ {*} \quad \text { when } \Phi \text { and } \Phi^ {\prime} \text { are bilinear }; \\ (c u _ {1}) ^ {*} = c ^ {J ^ {\prime}}. u _ {1} ^ {*} \quad \text { when } \Phi \text { and } \Phi^ {\prime} \text { are sesquilinear }. \end{array} \right.\tag{32}
\]

此外，若 \(E^{\prime\prime}\) 是第三个左 A-模，\(F^{\prime\prime}\) 是第三个右（相应地左）A-模，\(\Phi^{\prime\prime}\) 是 \(E^{\prime\prime} \times F^{\prime\prime}\) 上的一个双线性型（相应地关于 J 的半双线性型），且 \(u^{\prime}\) 是 \(E^{\prime}\) 到 \(E^{\prime\prime}\) 中的一个同态，其（左）伴随处处有定义，则有

\[
(u ^ {\prime} \circ u) ^ {*} = u ^ {*} \circ u ^ {\prime *}.\tag{33}
\]

特别地，若 \(u\) 是 E 到 E′ 上的一个同构，且伴随 \(u^{*}\) 与 \((u^{-1})^{*}\) 都处处有定义，则 \(u^{*}\) 是 F′ 到 F 上的一个同构，且我们有 \((u^{*})^{-1} = (u^{-1})^{*}\) 。右伴随有类似的性质。

命题 7. --- 在与上面相同的记号下，设 \(d_{\Phi}\) 是双射。于是每个 E 到 E′ 中的同态 \(u\) 都有一个处处有定义的左伴随，且我们有 \(u^{*} = (d_{\Phi})^{-1} \circ {}^{t}u \circ d_{\Phi'}\) 。

事实上，由于 \(d_{\Phi}\) 是双射，用本节开头的记号，有 \(\mathbf{F}_{1}^{\prime}=\mathbf{F}^{\prime}\) ，因而 \(u^{*}\) 处处有定义。另一方面，(30) 等价于

\[
\left\langle u (x), d _ {\Phi^ {\prime}} \left(y ^ {\prime}\right) \right\rangle = \left\langle x, \left(d _ {\Phi} \circ u ^ {*}\right) \left(y ^ {\prime}\right) \right\rangle \quad \left(x \in \mathrm{E}, y ^ {\prime} \in \mathrm{F} ^ {\prime}\right);
\]

或者 \(\langle d_{\Phi'}(y'), u(x) \rangle = \langle {}^t u(d_{\Phi'}(y')), x \rangle\) ；于是 \({}^t u(d_{\Phi'}(y')) = d_{\Phi}(u^*(y'))\) （对一切 \(y' \in F'\) ），从而 \({}^t u \circ d_{\Phi'} = d_{\Phi} \circ u^*\) ，因此得到所宣布的 \(u^*\) 的表达式。证毕。

注记. --- 当 \(s_{\Phi}\) 是双射时，每个 F 到 F′ 中的同态 \(\nu\) 都有一个处处有定义的右伴随，且我们有

\[
v ^ {*} = (s _ {\Phi}) ^ {- 1} \circ {} ^ {t} v \circ s _ {\Phi^ {\prime}}.\tag{34}
\]

命题 8. --- 在与上面相同的记号下，设 \(s_{\Phi}\) 与 \(d_{\Phi}\) 都是双射。设 \(u\) 与 \(v\) 分别是 E 到 \(\mathbf{E}'\) 上与 F 到 \(\mathbf{F}'\) 上的同构。于是，\(\Phi\) 是 \(\Phi'\) 关于 \(u\) 与 \(\nu\) 的原像（也就是说，有 \(\Phi(x,y)=\Phi'(u(x),\nu(y))\) （对一切 \(x\in\mathbf{E},y\in\mathbf{F}\) ））必须且只需我们有 \(u^{-1}=\nu^{*}\) 与 \(\nu^{-1}=u^{*}\) 。

事实上，\(\Phi'(u(x), v(y)) = \Phi(x, y)\) 也可写成 \(\Phi(x, u^*(v(y))) = \Phi(x, y)\) 。若它对一切 \(x \in E\) 与 \(y \in F\) 成立，则有 \(u^* \circ v = 1\) （由于 \(\Phi\) 非退化）。于是由 (33)，同样有 \(v^* \circ u = 1\) 。其逆是显然的。

推论. --- 设 A 是一个配备了反自同构 J 的环，E 是一个左 A-模，\(\Phi\) 是 E × E 上的一个（关于 J 的）半双线性型，其相伴映射都是双射，u 是 A-模 E 的一个自同构。u 保持 \(\Phi\) 不变（也就是说，对 E 中一切的 x, y 有 \(\Phi(u(x), u(y)) = \Phi(x, y)\) ）必须且只需 u 的两个伴随相等，且我们有 \(u^{*} = u^{-1}\) 。

这由命题 8 立即得出。

注记. --- 在命题 8 的推论的假设下，再设 A 是一个体且 E 在 A 上是有限维的。设 \(w\) 是 E 的一个自同态，\(w_{1}\) 与 \(w_{2}\) 是它的右伴随与左伴随。条件 \(ww_{1}=1\) 、\(ww_{2}=1\) 、\(w_{1}w=1\) 、\(w_{2}w=1\) 中的每一个都蕴涵 \(w\) 是 E 的一个保持 \(\Phi\) 不变的自同构，且 \(w_{1}=w_{2}\) 。

